\chapter{Energía}

\marginal{el agregado no se mueve, la composición sí}
La función Combustibles y Energía, Gobierno Federal más entidades, en bruto, pasa
de \textbf{1,361,059.8 a 1,373,266.1 millones de pesos}: 0.9\% nominal, es decir
una caída real de 3.2\%. El agregado apenas se mueve. Todo lo interesante está en
la composición.

\begin{table}[htbp]
\centering
\caption{Función Combustibles y Energía (mdp)}
\label{tab:C9_1}
\footnotesize
\setlength{\tabcolsep}{4pt}
\begin{tabular}{L{4.2cm}SSSSS}
\toprule
{Concepto} & {2024} & {2025} & {Dif. nominal} & {Dif. real} & {Var. real \%} \\
\midrule
{TOTAL función Combustibles y Energía (GF + entidades, bruto)} & \num{1361059.8} & \num{1373266.1} & {---} & {---} & {---} \\
\multicolumn{6}{l}{\itshape Gobierno Federal, por ramo} \\
{18 Energía} & \num{166789.3} & \num{137466.4} & \num{-29323.0} & \num{-36494.9} & \num{-21.0} \\
{23 Provisiones Salariales y Económicas} & \num{97281.3} & \num{104591.0} & \num{7309.7} & \num{3126.6} & \num{3.1} \\
{45 Comisión Reguladora de Energía} & \num{285.5} & \num{199.9} & \num{-85.7} & \num{-97.9} & \num{-32.9} \\
{46 Comisión Nacional de Hidrocarburos} & \num{247.6} & \num{173.3} & \num{-74.3} & \num{-84.9} & \num{-32.9} \\
\multicolumn{6}{l}{\itshape Entidades} \\
{TYY Pemex Consolidado} & \num{623985.1} & \num{611323.4} & \num{-12661.8} & \num{-39493.1} & \num{-6.1} \\
{TVV CFE Consolidado} & \num{472471.0} & \num{519512.3} & \num{47041.3} & \num{26725.0} & \num{5.4} \\
\bottomrule
\end{tabular}
\par\vspace{2pt}\footnotesize Fuente: elaboración propia del ITED con los analíticos del PEF aprobado 2024 y 2025. Las diferencias en millones de pesos se reportan dos veces, nominal y real; la real está en pesos de 2025 con el deflactor de 4.3 \% que declara el CGPE.
\end{table}


\section{Menos aportación fiscal, más ingreso propio}

La aportación del Gobierno Federal por el Ramo 18 cae de \textbf{167,736.2 a
138,307.4 millones de pesos}, 20.9\% real. Al mismo tiempo, y como mostró el
capítulo 3, el ingreso propio de las dos empresas sube con fuerza: Pemex de
769,805.6 a 860,868.2 y la Comisión Federal de Electricidad de 446,951.3 a
539,145.6.

\textbf{El paquete reduce el apoyo fiscal a las empresas energéticas y supone que
sus ingresos propios lo compensan.} Es la apuesta del capítulo y el paquete no
publica el supuesto de precios y volúmenes por empresa que la sostendría.

\begin{table}[htbp]
\centering
\caption{Pemex y CFE: por dónde entra y en qué se va (mdp)}
\label{tab:C9_2}
\small
\begin{tabular}{lSSSSS}
\toprule
{Concepto} & {2024} & {2025} & {Dif. nominal} & {Dif. real} & {Var. real \%} \\
\midrule
Pemex: gasto total (bruto) & \num{624805.6} & \num{612145.9} & \num{-12659.7} & \num{-39526.3} & \num{-6.1} \\
Pemex: pensiones (tipo de gasto 4) & \num{74120.5} & \num{83827.6} & \num{9707.1} & \num{6519.9} & \num{8.4} \\
Pemex: obra publica (capitulo 6000) & \num{237574.8} & \num{199216.9} & \num{-38357.9} & \num{-48573.6} & \num{-19.6} \\
Pemex: servicios personales (capitulo 1000) & \num{105084.2} & \num{114074.0} & \num{8989.8} & \num{4471.2} & \num{4.1} \\
Pemex: deuda publica (capitulo 9000) & \num{143341.2} & \num{147890.7} & \num{4549.5} & \num{-1614.2} & \num{-1.1} \\
CFE: gasto total (bruto) & \num{528532.7} & \num{583951.1} & \num{55418.4} & \num{32691.5} & \num{5.9} \\
CFE: pensiones (tipo de gasto 4) & \num{55961.3} & \num{64337.9} & \num{8376.6} & \num{5970.3} & \num{10.2} \\
CFE: obra publica (capitulo 6000) & \num{32405.2} & \num{39064.5} & \num{6659.3} & \num{5265.9} & \num{15.6} \\
CFE: servicios personales (capitulo 1000) & \num{74169.5} & \num{80447.5} & \num{6278.0} & \num{3088.7} & \num{4.0} \\
CFE: deuda publica (capitulo 9000) & \num{35152.0} & \num{38479.5} & \num{3327.5} & \num{1816.0} & \num{5.0} \\
Ramo 18 Energia (aportacion del Gobierno Federal) & \num{167736.2} & \num{138307.4} & \num{-29428.8} & \num{-36641.5} & \num{-20.9} \\
\bottomrule
\end{tabular}
\par\vspace{2pt}\footnotesize Fuente: elaboración propia del ITED con los analíticos del PEF aprobado 2024 y 2025. Las diferencias en millones de pesos se reportan dos veces, nominal y real; la real está en pesos de 2025 con el deflactor de 4.3\% que declara el CGPE.
\end{table}


\section{En qué se va}

Petróleos Mexicanos gasta 612,145.9 millones de pesos, 6.1\% menos en términos
reales. La caída se concentra en obra pública, que pasa de 237,574.8 a
\textbf{199,216.9 millones}, 19.6\% real. Mientras tanto, sus pensiones suben
8.4\% real y sus servicios personales 4.1\%, y su servicio de deuda se mantiene en
147,890.7.

La Comisión Federal de Electricidad hace lo contrario: gasta 5.9\% más en términos
reales, y \textbf{sube su obra pública 15.6\%}, de 32,405.2 a 39,064.5.

\marginal{una empresa invierte y la otra no}
El contraste es el hallazgo del capítulo. \textbf{Con menos apoyo fiscal, la
empresa eléctrica invierte más y la petrolera invierte menos.} En Pemex, la
combinación de obra que cae 19.6\% con pensiones que suben 8.4\% y deuda que no
cede describe una empresa cuyo presupuesto se lo llevan compromisos anteriores.

\section{Lo que falta}

El subsidio eléctrico no se puede aislar con las piezas disponibles, y las
estrategias programáticas por ramo, que explicarían el destino de la inversión de
cada empresa, no están en carpeta. Las dos ausencias se declaran.

\clearpage
